\begin{textsample}{POS Dim 2 – persona\_gemini – Score 20.00 – t050\_gemini.txt}  \label{ex:f2_pos_044}
An equitable and sustainable \textbf{future} is impossible without recognizing the deep interconnection between \textbf{climate} \textbf{justice} and social \textbf{justice}. \textbf{Climate} change \textbf{acts} as a \textbf{crisis} multiplier, exacerbating existing \textbf{inequalities}.
It disproportionately harms the very \textbf{communities} who have contributed the least to the \textbf{crisis}, including low-income countries, \textbf{people} of color, Indigenous Peoples, \textbf{women}, and those with disabilities.
These vulnerable \textbf{communities} are on the \textbf{frontlines}, \textbf{facing} devastating disasters, \textbf{resource} scarcity, and staggering \textbf{loss} and damage costs, which are projected to reach between 290 and 580 billion US dollars annually by 2030.
The \textbf{driver} of this \textbf{crisis} is clear : fossil \textbf{fuels} are responsible for over 75 % of all greenhouse gas emissions.
Addressing this requires a just phase-out of fossil \textbf{fuels} and a fundamental shift away from extractive \textbf{economies} towards people-centered models.
Polluters must be held accountable for the damage they have caused.
This is not a distant threat.
In the Middle East and North Africa region, for instance, Greenpeace has documented how escalating warming, \textbf{droughts}, and \textbf{floods} are worsening social \textbf{injustices}.
These \textbf{climate} \textbf{impacts} threaten the region ’s vital \textbf{ecosystems}, \textbf{people} 's \textbf{livelihoods}, and their cultural heritage.

% matched lemmas: act, climate, community, crisis, driver, drought, economy, ecosystem, face, flood, frontline, fuel, future, impact, inequality, injustice, justice, livelihood, loss, people, resource, woman
\end{textsample}
